\documentclass[10pt,a4paper]{amsart}
\usepackage[margin=19mm]{geometry}
\usepackage{amsmath,amssymb,mathtools}
\usepackage[scr=rsfs,cal=euler]{mathalfa}
\usepackage{xfrac}
\usepackage[unicode,hidelinks,bookmarks=false]{hyperref}
\newcommand{\ie}{i.e.\ }
\newcommand{\cf}{cf.\ }
\newcommand{\resp}{respectively\ }
\newcommand{\Subsection}{\subsection}
\providecommand{\nobreakdash}{\nobreak\hbox{-}}
\setlength{\emergencystretch}{2em}
\multlinegap0pt
\usepackage{mathtools}
\usepackage{dsfont}
\usepackage{xcolor}
\usepackage{tikz}
\usepackage{booktabs}
\usepackage[normalem]{ulem}
\usepackage{cancel}
\usepackage{float}
\usepackage{bm}
\usepackage{natbib}
\newtheorem{thm}{Theorem}
\usepackage{cleveref}
\crefname{thm}{Theorem}{Theorems}
\newcommand{\Aut}{\mathsf{Aut}}
\newcommand{\Z}{\mathbb{Z}}
\newcommand{\sm}{\smallsetminus}
\title{\centering\MakeUppercase{On Matrix Product Factorization of Cayley graphs}}
\author{Allen W. Herman, Bobby Miraftab}
\date{Source: arXiv:2512.17110v1 (CC BY 4.0)}
\begin{document}
\maketitle

\noindent\emph{Source and licence.} \centering\MakeUppercase{On Matrix Product Factorization of Cayley graphs}. Version arXiv:2512.17110v1 (CC BY 4.0), distributed by arXiv under the Creative Commons Attribution 4.0 International licence, http://creativecommons.org/licenses/by/4.0/, as recorded in the arXiv licence metadata for this exact version and shown on the abstract page https://arxiv.org/abs/2512.17110v1. Full licence notice: \texttt{licenses/arxiv-2512.17110-CC-BY-4.0.txt}.\par\smallskip

\noindent\textit{Editorial scope.} Counted unit: the article's thm thm:Z2n-to-D2n-factorable (source0.tex lines 1202--1232). The article's complete original proof of it is delivered with it (source0.tex lines 1234--1363, one \texttt{proof} environment of 130 lines). No mathematics is written, repaired or re-derived: the statement and the proof are the article's own lines, in the article's own notation, and no retained line is substituted or re-flowed.  Retained uncounted as the source-local context the proof needs: lines 324--332.  Nothing was omitted from any retained span, so the package carries no omission record.  No retained line contains a citation, so the wrapper carries no bibliography and no bibliography entry.  No retained line contains a figure environment, an image-inclusion command or drawing code, so no image is delivered and the package declares no asset dependency.  Editorial formatting only, in addition: an AMS \texttt{amsart} preamble that restates the article's own declarations and macros used by the retained lines, namely the environments \texttt{thm} on separate counters, together with the standard packages those lines need.  The retained environments print in this document as Theorem 1 in order of appearance, whereas the article prints the same items with the numbering of its own full document.  The complete native source of the article as distributed in the arXiv e-print for this exact version is bundled unchanged as \texttt{sources/191290/source0.tex}.
\begin{thm}\label{thm:main_1}
Let $G$ be a finite group and let $S,T,U \subseteq G\sm \{e\}$ be symmetric sets. 
Then the following statements are equivalent:
\begin{enumerate}
    \item $(S,T,U)$ is factorable in $G$.
    \item $U=ST$ and moreover for every $u\in U $, there exists a unique pair $(s,t)\in S\times T$ such that $u=st$.
    \item $(\sum_{s\in S} s)(\sum_{t\in T}t )=\sum_{u\in U}u$.
\end{enumerate}
\end{thm}

\begin{thm}\label{thm:Z2n-to-D2n-factorable}
Let $n\ge 3$ be odd. 
Then the following holds:
\begin{enumerate}
\item If $(S,T,U)$ is factorable in $\Z_{2n}$,
then $(\widehat S,\widehat T,\widehat U)$ is factorable in $D_{2n}$.
Moreover, if $(S,T,U)$ is equivalent to $(A,B,C)$ in $\Z_{2n}$, then
%via an automorphism
%$x\mapsto ux$ with $u\in\Z_{2n}^\times$
$(\widehat S,\widehat T,\widehat U)$ is equivalent to $(\widehat A,\widehat B,\widehat C)$
in $D_{2n}$.
%via the automorphism $f\in\Aut(D_{2n})$ defined by $f(r)=r^{\,u\bmod n}$, $f(s)=s$.
\item If $(\widetilde S,\widetilde T,\widetilde U)$ is factorable in $D_{2n}$, then
$(S,T,U)$ is factorable in $\Z_{2n}$. 
If, in addition,
$|\widetilde S|=k$, $|\widetilde T|=\ell$ satisfy
\[
\gcd\!\Bigl(n,\tfrac{k+1}{2}\Bigr)=\gcd\!\Bigl(n,\tfrac{\ell+1}{2}\Bigr)=1,
\]
then equivalence in $D_{2n}$ pulls back to equivalence in $\Z_{2n}$ in the sense that, whenever
$(\widetilde S,\widetilde T,\widetilde U)$ is equivalent to $(\widetilde A,\widetilde B,\widetilde C)$ in
$D_{2n}$, the triples
\[
\bigl(\varphi^{-1}\psi^{-1}(\widetilde S),\ \varphi^{-1}\psi^{-1}(\widetilde T),\ \varphi^{-1}\psi^{-1}(\widetilde U)\bigr)
\quad\text{and}\quad
(A,B,C)
\]
are equivalent in $\Z_{2n}$.
%(again via $x\mapsto ux$ for some $u\in\Z_{2n}^\times$).
\end{enumerate}
\end{thm}

\begin{proof}
By~\Cref{thm:main_1}, we know that $(S,T,U)$ is factorable in $\Z_{2n}$ if and only if 
the “unique-sum” counts satisfy
\[
N_{S,T}(x)=
\begin{cases}
1,&x\in U,\\
0,&x\notin U,
\end{cases}
\qquad
N_{S,T}(x)\coloneqq |\{(s,t)\in S\times T\mid \ s+t=x\}|.
\]
Next we set the following sets
$S^*\coloneqq \varphi(S)$, $T^*\coloneqq \varphi(T)$, $U^*\coloneqq \varphi(U)\subseteq\Z_2\times\Z_n$.
We shall show that the same uniqueness holds for componentwise opearation in $\Z_2\times\Z_n$.
We note that every $(\alpha,\beta)\in U^*$ has a \emph{unique} representation
$(\alpha,\beta)=(i,k)+(i',\ell)$ with $(i,k)\in S^*$, $(i',\ell)\in T^*$, as we desired.\\
Let $\widehat S=\psi(S^*)$, $\widehat T=\psi(T^*)$, $\widehat U=\psi(U^*)$.
We show that the “unique-product” condition holds in $D_{2n}$.
We pick up an element $x\in D_{2n}\sm \{e\}$.
Then we have the following cases:
\begin{itemize}
\item If $x=r^j$ is a rotation, write $(0,j)=(i,k)+(i',\ell)$ in $\Z_2\times\Z_n$ with
$(i,k)\in S^*$, $(i',\ell)\in T^*$. If $i=i'=0$ then
$\psi(i,k)\psi(i',\ell)=r^{k}r^{\ell}=r^j$. If $i=i'=1$ then
$\psi(i,k)\psi(i',\ell)=(sr^{k})(sr^{\ell})=r^{\ell-k}$; by symmetry of $S^*,T^*$
we can replace $k$ by $-k$ to get $r^{(-k)+\ell}=r^j$. 
\item If $x=sr^j$ is a reflection, write $(1,j)=(i,k)+(i',\ell)$. If $(i,i')=(0,1)$ then
$\psi(i,k)\psi(i',\ell)=r^{k}(sr^{\ell})=sr^{k+\ell}=sr^j$. 
If $(i,i')=(1,0)$ then
$(sr^{k})r^{\ell}=sr^{k+\ell}=sr^j$.
\end{itemize}

Thus every non-identity element of $D_{2n}$ lies in $\widehat S\,\widehat T$.
Next we shall show the uniqueness.
Suppose $x=\widehat s_1\widehat t_1=\widehat s_2\widehat t_2$ with
$\widehat s_i\in\widehat S$, $\widehat t_{\ell}\in\widehat T$ for $\ell=1,2$.
Apply $\psi^{-1}$ to obtain two representations of either $(0,j)$ or $(1,j)$ in
$\Z_2\times\Z_n$ as $(i,k)+(i',\ell)$. Uniqueness in $\Z_2\times\Z_n$ forces
$(i,k)=(i',k')$ and $(i',\ell)=(i'',\ell')$, hence $\widehat s_1=\widehat s_2$ and
$\widehat t_1=\widehat t_2$. Therefore the product representation in $D_{2n}$ is unique.
It follows from~\Cref{thm:main_1} that the uniqueness of representations is equivalent to
$A(D_{2n};\widehat U)=A(D_{2n};\widehat S)\,A(D_{2n};\widehat T)$, as required.

Next we show the equivalence under automorphisms.
More precisely we show
if $(S,T,U)$ is equivalent to $(A,B,C)$ in $\Z_{2n}$ via an automorphism
$x\mapsto ux$ with $u\in\Z_{2n}^\times$, then 
$(\widehat S,\widehat T,\widehat U)$ is equivalent to $(\widehat A,\widehat B,\widehat C)$
in $D_{2n}$ via the automorphism $f\in\Aut(D_{2n})$ defined by $f(r)=r^{\,u\bmod n}$, $f(s)=s$.

If $(A,B,C)=(uS,uT,uU)$ with $u\in\Z_{2n}^\times$, then
$S^*\mapsto f(S^*)$, $T^*\mapsto f(T^*)$, $U^*\mapsto f(U^*)$ under
$f(i,j)=(i,\,u j)$ on $\Z_2\times\Z_n$. 
Next we consider the following automorphism of $D_{2n}$
\[
\begin{array}{rcl}
f_{u,0} \colon D_{2n} &\longrightarrow& D_{2n} \\[4pt]
r &\longmapsto& r^{\,u}, \\[2pt]
s &\longmapsto& s
\end{array}
\]
Then it is not hard to see that 
$\psi\circ f=f_{u,0}\circ \psi$
 which implies $(\widehat S,\widehat T,\widehat U)\mapsto
\bigl(f_{u,0}(\widehat S),\,f_{u,0}(\widehat T),\,f_{u,0}(\widehat U)\bigr)$,
i.e., equivalence is preserved. 

Let $(\widetilde S,\widetilde T,\widetilde U)$ is factorable in $D_{2n}$.
Then we set
$S^*=\psi^{-1}(\widetilde S)$, $T^*=\psi^{-1}(\widetilde T)$, $U^*=\psi^{-1}(\widetilde U)$
in $\Z_2\times\Z_n$ and then $S=\varphi^{-1}(S^*)$, etc.
The same two product patterns as above show that every nonzero $(\alpha,\beta)$
has a unique decomposition as a sum from $S^*,T^*$.
By applying $\varphi^{-1}$, we derive the uniqueness in $\Z_{2n}$, hence $A(\Z_{2n};U)=A(\Z_{2n};S)A(\Z_{2n};T)$.

We need to show that
$(\widetilde S,\widetilde T,\widetilde U)$ and $(\widetilde A,\widetilde B,\widetilde C)$
are equivalent if there exists $f\in\Aut(D_{2n})$ with
$(\widetilde A,\widetilde B,\widetilde C)=(f(\widetilde S),f(\widetilde T),f(\widetilde U))$.
For odd $n$, every automorphism has the form
\[
f_{u,v}:\quad r\mapsto r^{\,u},\qquad s\mapsto s\,r^{\,v},
\qquad (u\in\Z_n^{\times},\ v\in\Z_n).
\]
Conjugating by $\psi$ shows the induced action on $\Z_2\times\Z_n$ is the affine map
\[
F_{u,v}:\ (i,j)\longmapsto\bigl(i,\;u\,j+i\,v\bigr).
\]
Thus if $(\widetilde A,\widetilde B,\widetilde C)=\bigl(f_{u,v}(\widetilde S),f_{u,v}(\widetilde T),
f_{u,v}(\widetilde U)\bigr)$, then with $S^*=\psi^{-1}(\widetilde S)$, etc.,
\[
A^*=F_{u,v}(S^*),\qquad B^*=F_{u,v}(T^*),\qquad C^*=F_{u,v}(U^*).
\]
Applying $\varphi^{-1}$ gives the pulled–back triple $(A,B,C)\subseteq\Z_{2n}$.


Let us fix $k=|\widetilde S|$, $\ell=|\widetilde T|$, and let
$s\coloneqq |\widetilde S\cap s\langle r\rangle|$, $t\coloneqq |\widetilde T\cap s\langle r\rangle|$
be the counts of reflections in the two factors. 
It is not hard to see that 
\[
s=\frac{k+\varepsilon}{2},\qquad t=\frac{\ell-\varepsilon}{2}
\quad\text{for some }\varepsilon\in\{+1,-1\};
\]
in particular, one of $\frac{k+1}{2},\frac{k-1}{2}$ equals $s$, and one of
$\frac{\ell+1}{2},\frac{\ell-1}{2}$ equals $t$.

Under $F_{u,v}$, the rotation indices (the $i=0$ slice of $S^*$) are multiplied by $u$,
while the reflection indices (the $i=1$ slice) are sent to $u\,M_S+v$, where
$M_S\coloneqq \{\,j:\ (1,j)\in S^*\,\}\subseteq\Z_n$.
If the pulled–back equivalence on $\Z_{2n}$ is to be a pure multiplier $x\mapsto u'x$,
then necessarily $v\equiv 0\pmod n$; otherwise the reflection index set would be
translated by a nonzero $v$. To force $v\equiv 0$, observe that
\[
u\,M_S+v\ =\ M_A\quad\Longrightarrow\quad M_S+v\ =\ u^{-1}M_A,
\]
so $M_S$ is invariant under translation by $v$. Hence $M_S$ is a union of cosets of
the subgroup $\langle v\rangle\le\Z_n$, and therefore $|M_S|$ is a multiple of
$n/\gcd(n,v)$. But $|M_S|=s$ equals either $\tfrac{k+1}{2}$ or $\tfrac{k-1}{2}$.
If, in particular, $\gcd\!\bigl(n,\tfrac{k+1}{2}\bigr)=1$, then the only way a nonempty
subset of size $s$ can be a union of cosets of $\langle v\rangle$ is $v\equiv 0\pmod n$.
Applying the same argument to $T$ and the hypothesis
$\gcd\!\bigl(n,\tfrac{\ell+1}{2}\bigr)=1$ yields $v\equiv 0\pmod n$ as well.

Consequently the affine map reduces to $F_{u,0}(i,j)=(i,uj)$, which after $\varphi^{-1}$
becomes the multiplier $x\mapsto u'x$ on $\Z_{2n}$ (with $u'\equiv u\pmod n$, $u'\equiv 1\pmod 2$).
Thus the pulled–back triples $(A,B,C)$ and $(S,T,U)$ are equivalent in $\Z_{2n}$ via
an automorphism $x\mapsto u'x$.
\end{proof}

\end{document}
